% Drafted 2026-06-25. Durable (NOT pending-144): the instrument-validation outcome is
% verified now. Sources: results/ab_gate/amended_cceb325_result.json,
% docs/ab-gate-findings-2026-06-24.md, docs/amended-gate-2026-06-24.md.
% INTEGRITY: the +15.3pp A/B resolve delta is "the fixed instrument is net-positive",
% NOT a thesis result. Never state the original gate passed.

% ── 5.0  Instrument-Validation Result ────────────────────────────────────
% PURPOSE: Report the verified outcome of the pre-registered A/B instrument health gate.
% ─────────────────────────────────────────────────────────────────────────

\section{Instrument-Validation Result}
\label{sec:instrument-results}

The instrument-health procedure of Section~\ref{sec:instrument} produced a verified outcome
that gates the fixed-code rerun \emph{before} the matrix is trusted, and unlike the H1--H5
contrasts these numbers do not depend on that rerun --- they are reported here as a durable
methodological result.

\paragraph{From STOP to GO.}
On the original pre-registered gate the instrument \textbf{STOPped}: the fixed-mode arm
showed 82 edit path/index failures and an edit-failure ratio of $0.303$, of which the
dominant component was the path-normalisation bug --- 77 of 78 ``security'' rejections were
false (the same file under two normalisation states). After the hotfix (\texttt{cceb325}),
the instrument-attributable failures fall from 79 to \textbf{0}. The amended gate, evaluated
on fresh post-hotfix data, returns \textbf{GO}: all of its criteria pass
(Table~\ref{tab:instrument}). Instrument-attributable failures are zero in both arms; the
model-quality failure rate in the fixed arm ($0.152$) is below the legacy arm ($0.183$), so
the fix does not regress diff quality; and token inflation is below $1.0\times$ on both the
prompt ($0.986\times$) and total ($0.992\times$) medians, comfortably under the $1.5\times$
ceiling. Completeness, pairing, and provenance (commit, configuration hash, and row-level
tool mode) all hold.

\begin{table}[htbp]
  \centering
  \caption{Amended instrument-health gate on fresh post-hotfix data. Instrument-attributable
  failures are zero in both arms; the fixed arm's model-quality rate does not exceed the
  legacy arm's; token inflation is below $1.0\times$. Verdict: GO.}
  \label{tab:instrument}
  \begin{tabular}{lrr}
    \toprule
    & Legacy (pre-fix tooling) & Fixed (post-fix tooling) \\
    \midrule
    Edit calls                       & 339      & 532      \\
    Instrument-attributable failures & 0        & 0        \\
    Model-quality failures           & 62       & 81       \\
    Model-quality rate               & 0.183    & 0.152    \\
    Prompt-token median              & 177{,}322 & 174{,}900 \\
    Total-token median               & 180{,}352 & 178{,}825 \\
    \midrule
    Prompt inflation (fixed/legacy)  & \multicolumn{2}{c}{$0.986\times$} \\
    Total inflation (fixed/legacy)   & \multicolumn{2}{c}{$0.992\times$} \\
    Verdict                          & \multicolumn{2}{c}{\textbf{GO}} \\
    \bottomrule
  \end{tabular}
  \par\smallskip
  {\footnotesize Source: \texttt{results/ab\_gate/amended\_cceb325\_result.json}.}
\end{table}

\paragraph{The fixed instrument is net-positive (informational).}
Beyond passing the health criteria, the corrected tooling resolved \emph{more} tasks than
the legacy tooling on the A/B's two sequences: a resolution rate of $0.582$ fixed versus
$0.429$ legacy ($+15.3$ percentage points). This is reported as informational evidence that
the instrument fix is net-positive --- it confirms that the read/edit corrections help the
agent rather than merely changing failure bookkeeping. It is \emph{not} a thesis result: it
is measured on two sequences and three policies under the instrument check, not on the
144-run matrix, and the headline between-policy contrasts are reported from the trustworthy
rerun in the sections that follow.
